\section{Data Construction}
\label{sec:data}

Existing privacy benchmarks for LLMs are predominantly organized
around \emph{explicit personally identifiable information} (PII),
e.g., names, emails, addresses, and organization
names~\citep{pilan2022text,lukas2023analyzing}.%
\footnote{Citations to be verified against the camera-ready bibliography.}
Yet in real interactions, users routinely disclose information that
is sensitive but \emph{not identifying}---family relationships,
health conditions, financial pressure, whereabouts, third-party
identities---much of which is \emph{unnecessary} for the task at
hand. Following the contextual-integrity view of
privacy~\citep{nissenbaum2010privacy,mireshghallah2024can,ngong2025contextual},
we frame the rewriter's target not as PII redaction but as the
suppression of \textbf{task-non-essential sensitive disclosures},
while preserving information that is \emph{contextually essential} to
answering the user's request.

Concretely, every sample in our benchmark comes with two disjoint
attribute sets that directly instantiate the symbols used in
\S\ref{sec:reward}: a set of \textbf{essential} spans
$\mathcal{E}=\{e_i\}_{i=1}^{n}$ that must be preserved, and a set of
\textbf{non-essential} spans $\mathcal{N}=\{c_j\}_{j=1}^{K}$ that
should be removed or generalized. To populate
$(\mathcal{E},\mathcal{N})$ with both breadth and label reliability,
we build two complementary subsets: (i) a \textbf{general-domain}
subset constructed via a reverse-synthesis pipeline
(\S\ref{sec:data:reverse}), augmented with a small but high-purity
slice of \emph{natural} single-turn user prompts to anchor the
distribution to real usage; and (ii) a \textbf{medical} challenge
set that stresses the over-removal failure mode in a domain where
sensitive information is often the answer itself
(\S\ref{sec:data:medical}). Detailed prompts, value banks, and
per-domain pools are deferred to the appendix.

\subsection{General-Domain Samples}
\label{sec:data:reverse}

The general-domain subset is dominated by reverse-synthesized prompts
and anchored by a small slice of real WildChat queries; both routes
are labeled under the same $(\mathcal{E},\mathcal{N})$ schema and
the same heterogeneous-judge protocol described below.

\paragraph{Reverse-synthesis route.}
We adopt a \textbf{reverse-synthesis} design: we first sample
a structured \emph{semantic skeleton}---domain, task, subject, and
the partition of sensitive attributes into $\mathcal{E}$ and
$\mathcal{N}$---and only then ask a generator to embed this skeleton
into a fluent user prompt. An independent judge model re-extracts
$(\mathcal{E},\mathcal{N})$ from the surface text and is required to
agree with the sampler before a sample is admitted. This closed loop
produces labels that are \emph{generated by construction} rather than
inferred after the fact.

\paragraph{Structured sampler.}
Following the taxonomy of \citet{ngong2025contextual},%
\footnote{Citation to be verified.}
the sampler covers \textbf{11 privacy-relevant domains} (e.g., health,
finance, employment), \textbf{20 task types} (e.g., summarization,
translation, code generation), and \textbf{30 sensitive attribute
types} (e.g., family history, medication). To inject a linguistically
plausible prior over which attributes can be \emph{essential}, we
attach to each domain a hand-curated \emph{essential candidate pool}
used as a \textbf{soft} prior---essential attributes are preferably,
but not exclusively, drawn from this pool. To prevent stylistic
collapse, we further sample three orthogonal axes: the
\emph{speaker subject}, the \emph{embedding style}, and the
\emph{emotional tone}.

\paragraph{Generator and consistency judge.}
A generator LLM is conditioned on the sampled skeleton and instructed
to embed every non-essential attribute naturally while making each
essential attribute \emph{semantically load-bearing}, i.e., its
removal should make the request unanswerable at the same quality. To
avoid co-adaptation between the generator and the labeler, we
deliberately use \textbf{heterogeneous models}: Claude Sonnet~4.6%
\footnote{Citation to be verified.}
as the generator and GPT-4o%
\footnote{Citation to be verified.}
as the judge. The judge sees only the surface prompt and re-extracts
a triple ($\emph{primary context}$, $\hat{\mathcal{E}}$,
$\hat{\mathcal{N}}$). We then match $\hat{\mathcal{E}}$ against the
sampler's $\mathcal{E}$ (and likewise for $\mathcal{N}$) using a
two-stage matcher: first, a stop-word-filtered \textbf{token-overlap}
coverage with threshold $0.5$; if it fails, we fall back to
\textbf{embedding cosine similarity} with threshold $0.65$ to absorb
Only samples whose extracted labels agree with
the sampler at both ends of the partition are kept. We provide the
full prompts, thresholds, and failure-mode statistics in the
appendix.

\paragraph{Natural-dialogue anchor.}
Reverse synthesis gives labels that are correct \emph{by
construction}, but at the cost of a controlled distribution that may
drift from how users actually phrase sensitive requests. To verify
that the same $(\mathcal{E},\mathcal{N})$ protocol applies to
organic prompts, we curate a natural anchor set from
WildChat~\citep{zhao2024wildchat}%
\footnote{Citation to be verified.}
through a four-stage funnel. Starting from
\textbf{$\sim$352k} raw conversations, we first apply a
deterministic structural filter (single-turn, English, non-toxic,
non-redacted, $25$--$250$ words; this length range is where the
eliciting utterance---rather than a pasted document---fixes the
contextual norm), leaving \textbf{$\sim$50k} candidate prompts.
A lightweight rater (GPT-4o-mini) then scans this pool and retains
only prompts that exhibit a clear contextual-privacy violation,
yielding \textbf{656} candidates. Finally, GPT-4o re-extracts $(\mathcal{E},\mathcal{N})$, and we admit only those whose extracted
$\mathcal{N}$ is non-empty, leaving \textbf{345 anchored real-user
samples}. This set acts as an
out-of-distribution probe for rewriters trained on the synthesized
portion.

\subsection{Medical Challenge Set}\label{sec:data:medical}

Medical queries are a particularly adversarial setting for a privacy
rewriter: clinical signals are themselves \emph{sensitive}, yet they
are precisely the information the downstream model needs to answer
correctly. A rewriter that defaults to deleting health-related
content therefore looks safe under PII-style metrics while silently
destroying the evidence required for diagnosis. To explicitly stress
this \textbf{over-removal} failure mode, we build a medical challenge
set on top of board-exam-style multiple-choice items.%
\footnote{Built on top of medical-o1-style data; citation to be verified.}

\paragraph{Essential extraction.}
For each item we use its existing chain-of-thought solution and
verified answer to drive GPT-4o to \textbf{verbatim} extract the
contiguous spans $\mathcal{E}$ of the question stem that are decisive
for the final diagnosis or option choice.

\paragraph{Controlled disclosure injection.}
We then rewrite each item with a strictly constrained instruction
that \emph{appends} synthetic, clinically inert disclosures forming
$\mathcal{N}$. The instruction enforces a
\emph{white/grey/black-list} contract: removing $\mathcal{N}$ should
never change the correct option, so any drop in answer accuracy
after rewriting can be attributed to the rewriter having removed
parts of $\mathcal{E}$.
